\begin{remark*}[Exposition]
A mathematical space is an arena in which a particular kind of mathematics can
be performed.  At the base there is a set of objects.  Added structure then
specifies the ground rules: which objects can be compared, which subsets are
distinguished, which operations are allowed, which limits make sense, or which
quantities can be measured.  The same underlying set can support many different
spaces because different structures make different questions meaningful.
\end{remark*}

\begin{remark*}[Structural Pattern]
The general pattern is
\[
  \text{space} = \text{underlying set} + \text{chosen structure}.
\]
A metric space chooses a distance function.  A topological space chooses a
collection of open sets.  A measurable space chooses a $\sigma$-algebra of
measurable sets, and a measure space adds a measure to that measurable
structure.  Algebraic spaces choose operations and laws.  Each choice controls
which statements are legitimate inside that setting.
\end{remark*}

\begin{figure}[htbp]
\centering
\begin{tikzpicture}[
  font=\small,
  box/.style={
    draw=black!45,
    rounded corners=2pt,
    align=center,
    inner xsep=8pt,
    inner ysep=6pt,
    minimum height=1cm
  },
  base/.style={box, fill=gray!8, minimum width=11.2cm},
  branch/.style={box, fill=blue!6, minimum width=3.2cm},
  refine/.style={box, fill=green!6, minimum width=3.2cm},
  measure/.style={box, fill=orange!8, minimum width=3.2cm},
  algebra/.style={box, fill=violet!7, minimum width=3.2cm},
  arrow/.style={-Latex, thick, draw=black!55}
]
\node[base] (set) at (0,0) {Underlying set \(X\)};
\node[base] (space) at (0,-1.45)
  {Mathematical space: \(X\) with selected structure};

\node[branch] (top) at (-3.9,-3.15)
  {Topological space\\open sets};
\node[refine] (metric) at (-3.9,-4.75)
  {Metric space\\distance \(d\)};
\node[refine] (normed) at (-3.9,-6.35)
  {Normed vector space\\linear structure + norm};

\node[measure] (measurable) at (0,-3.15)
  {Measurable space\\\(\sigma\)-algebra};
\node[measure] (measure) at (0,-4.75)
  {Measure space\\measure \(\mu\)};

\node[algebra] (algebraic) at (3.9,-3.15)
  {Algebraic structure\\operations + laws};

\draw[arrow] (set) -- (space);
\draw[arrow] (space) -- (top);
\draw[arrow] (top) -- (metric);
\draw[arrow] (metric) -- (normed);

\draw[arrow] (space) -- (measurable);
\draw[arrow] (measurable) -- (measure);

\draw[arrow] (space) -- (algebraic);

\draw[arrow, dashed] (metric.east) to[bend left=14] (top.east);
\end{tikzpicture}

\caption{Mathematical spaces are built by adding structure to an underlying
set.  Some structures refine others: a metric induces a topology, while a
measure requires measurable sets.}
\label{fig:nested-mathematical-spaces}
\end{figure}

\begin{remark*}[Interpretation]
The nesting in Figure~\ref{fig:nested-mathematical-spaces} should be read as
an increase in available structure, not as a claim that every subject lies on a
single chain.  Metric spaces sit naturally over topology because every metric
determines open balls and hence open sets.  Measure spaces follow a different
branch: they organize size rather than nearness.  The common feature is the
same: mathematical work begins only after the arena and its rules have been
specified.
\end{remark*}

\begin{figure}[htbp]
\centering
\begin{tikzpicture}[
  font=\small,
  space/.style={
    draw=#1!70,
    fill=#1!10,
    rounded corners=5pt,
    line width=0.8pt,
    align=center,
    minimum width=3.1cm,
    minimum height=1.05cm,
    inner sep=4pt
  },
  smallspace/.style={
    draw=#1!70,
    fill=#1!10,
    rounded corners=5pt,
    line width=0.8pt,
    align=center,
    minimum width=2.55cm,
    minimum height=0.9cm,
    inner sep=4pt,
    font=\scriptsize
  },
  note/.style={
    align=center,
    font=\scriptsize,
    text=black!70
  },
  arrow/.style={-Latex, thick, draw=black!55},
  softarrow/.style={-Latex, thick, dashed, draw=black!45}
]

\node[space=gray, minimum width=4.0cm] (sets) at (0,3.0)
  {\textbf{Sets}\\[-1pt]\scriptsize objects only};

\node[space=teal] (families) at (-3.55,1.45)
  {\textbf{Families of subsets}\\[-1pt]\scriptsize chosen subsets of \(X\)};
\node[space=blue] (topology) at (-3.55,-0.05)
  {\textbf{Topological spaces}\\[-1pt]\scriptsize open sets};
\node[smallspace=cyan] (metric) at (-3.55,-1.55)
  {\textbf{Metric spaces}\\[-1pt] distance \(d\)};

\node[space=green] (vector) at (1.25,1.45)
  {\textbf{Vector spaces}\\[-1pt]\scriptsize addition and scalars};
\node[smallspace=orange] (normed) at (-0.25,-0.05)
  {\textbf{Normed vector}\\\textbf{spaces}\\[-1pt] length};
\node[smallspace=purple] (inner) at (3.05,-0.05)
  {\textbf{Inner product}\\\textbf{spaces}\\[-1pt] angles and length};
\node[smallspace=red] (banach) at (-0.25,-1.55)
  {\textbf{Banach spaces}\\[-1pt] complete\\normed};
\node[smallspace=violet] (hilbert) at (3.05,-1.55)
  {\textbf{Hilbert spaces}\\[-1pt] complete\\inner product};

\node[space=brown] (measurable) at (5.8,1.45)
  {\textbf{Measurable spaces}\\[-1pt]\scriptsize \(\sigma\)-algebra};
\node[smallspace=red] (measure) at (5.8,-0.05)
  {\textbf{Measure spaces}\\[-1pt] size \(\mu\)};

\draw[arrow] (sets.south west) to[bend right=8] (families.north);
\draw[arrow] (families) -- (topology);
\draw[arrow] (sets.south) to[bend left=4] (vector.north);
\draw[arrow] (sets.south east) to[bend left=10] (measurable.north);

\draw[softarrow] (metric.north) -- (topology.south);
\draw[arrow] (vector) -- (normed);
\draw[arrow] (vector) -- (inner);
\draw[arrow] (normed) -- (banach);
\draw[arrow] (inner) -- (hilbert);
\draw[softarrow] (inner.south west) to[bend left=12] (normed.south east);
\draw[arrow] (measurable) -- (measure);

\node[note, text width=11.6cm] at (1.0,-2.85)
  {The branches record added structure.  A topology selects open sets, a
  metric supplies distance and induces open sets, linear spaces add operations,
  and measure spaces add measurable sets together with size.};

\end{tikzpicture}

\caption{One common route from sets to structured spaces.  A topology begins
with selected open subsets of a set.  A metric gives a distance and thereby
generates open balls, so metric spaces naturally carry topologies.  Linear,
normed, inner product, Banach, and Hilbert spaces add further compatible
structure.}
\label{fig:sets-to-spaces-hierarchy}
\end{figure}

\begin{remark*}[Exposition]
The hierarchy in Figure~\ref{fig:sets-to-spaces-hierarchy} is a guide to
structure, not a universal taxonomy.  Some spaces are formed by adding data to
a set.  Others are formed by adding compatibility requirements to an already
structured space.  The following table records the construction pattern that
will be used throughout this volume.
\end{remark*}

\begin{remark*}[Structure Table]
\begin{center}
\small
\renewcommand{\arraystretch}{1.2}
\begin{tabularx}{\linewidth}{>{\raggedright\arraybackslash}p{0.22\linewidth}
                            >{\raggedright\arraybackslash}p{0.24\linewidth}
                            >{\raggedright\arraybackslash}X}
\toprule
\textbf{Space} & \textbf{Comes from} & \textbf{Additional structure or rules} \\
\midrule
Set & Primitive arena & A collection of objects; no nearness, size, or
operation is specified. \\
Family of subsets & Set \(X\) & A selected collection
\(\mathcal{A}\subseteq\mathcal{P}(X)\). \\
Topological space & Set \(X\) & A topology \(\mathcal{T}\): open sets
containing \(\varnothing\) and \(X\), closed under arbitrary unions and finite
intersections. \\
Metric space & Set \(X\) & A metric \(d:X\times X\to[0,\infty)\) satisfying
identity of indiscernibles, symmetry, and the triangle inequality; the metric
induces open balls and hence a topology. \\
Vector space & Set \(V\) & Addition and scalar multiplication satisfying the
vector space axioms over a chosen field. \\
Normed vector space & Vector space & A norm \(\lVert\cdot\rVert\) measuring vector
length, compatible with scalar multiplication and addition. \\
Inner product space & Vector space & An inner product \(\langle\cdot,\cdot\rangle\)
measuring angles and lengths; it induces a norm. \\
Banach space & Normed vector space & Completeness: every Cauchy sequence
converges in the norm. \\
Hilbert space & Inner product space & Completeness with respect to the norm
induced by the inner product. \\
Measurable space & Set \(X\) & A \(\sigma\)-algebra \(\mathcal{M}\) of subsets
regarded as measurable. \\
Measure space & Measurable space & A measure \(\mu:\mathcal{M}\to[0,\infty]\)
assigning size and satisfying countable additivity. \\
\bottomrule
\end{tabularx}
\end{center}
\end{remark*}
